%% =================================================================
%% Beber, Lamon, Saveriano, Fontanelli, Palopoli.
%% "Autonomous Robotic Tissue Palpation and Abnormalities
%%  Characterisation via Ergodic Exploration."
%% IEEE Robotics and Automation Letters, preprint accepted Feb. 2026.
%% DOI: 10.1109/LRA.2026.3673907
%%
%% Reconstruction of page 3 as study notes.
%% Equation-dense: numbered equations (1), (2), (3); several
%% unnumbered display equations (GP, SE kernel, coverage density).
%% =================================================================

\documentclass[11pt]{article}

\usepackage[margin=1in]{geometry}
\usepackage{amsmath, amssymb}
\usepackage{mathrsfs}
\usepackage{bm}
\usepackage{xcolor}
\usepackage{fancyhdr}
\usepackage{url}

\pagestyle{fancy}
\fancyhf{}
\lhead{\small Beber et al.: Ergodic Palpation for Tissue Mapping}
\rhead{\small IEEE RA-L Preprint, Accepted Feb.\ 2026}
\cfoot{\small 3 $\to$ \arabic{page}}
\renewcommand{\headrulewidth}{0.4pt}

\setcounter{page}{1}
\setcounter{equation}{0}   % pages 1--2 had no equations; first numbered eq on p.3 is (1)
\setlength{\parskip}{4pt}

\begin{document}

% ---------- Preprint notice strip ----------
\begin{center}
{\footnotesize\itshape
This article has been accepted for publication in IEEE Robotics and
Automation Letters. This is the author's version which has not been
fully edited and content may change prior to final publication.
Citation information: DOI 10.1109/LRA.2026.3673907\par}
\end{center}

\vspace{0.4em}
\hrule
\vspace{0.6em}

\noindent
{\footnotesize\scshape Beber et al.: Ergodic Palpation for Tissue Mapping \hfill 3}

\vspace{1.2em}

% ---------- Continuation of Methodology from page 2 ----------
\noindent
spatial gradients. Although viscoelastic modelling supports robust
estimation, lesion search and characterisation in this work are driven
solely by elasticity. Our framework is structured as a closed-loop
search algorithm comprising four main components: A. a target
distribution over tissue stiffness, which is updated at every
iteration; B. an Expected Information Density (EID) map to prioritise
sampling locations; C. a trajectory planner using ergodic control,
guided by the EID; D. a data acquisition executed during trajectory
execution, which enables updates in the target distribution. These
steps are executed cyclically, updating the target distribution (E.)
until a predefined stopping criterion is met (F.). Each component is
designed to be modular and independent, allowing flexible integration
of alternative models or control strategies.

\vspace{0.6em}

% ---------- A. Target Distribution ----------
\noindent\textit{A.~Target Distribution}

\vspace{0.2em}

\noindent
We use GPR to construct a continuous map of the stiffness distribution
across the palpated surface. The robot end-effector positions serve as
the input domain, and the estimated elasticity values form the output.
After each regression update, the mean function $\mu(x)$ estimates the
stiffness at unobserved locations, while the variance function
$\sigma^{2}(x)$ quantifies the model uncertainty. As expected with
GPR, uncertainty increases with distance from the observed samples.
The squared exponential kernel ensures smooth, differentiable
estimates, which are critical for downstream gradient computation.

\vspace{0.3em}

\noindent\textit{\hspace{1em}1)~Gaussian Process Regression:} is a
non-parametric Bayesian approach for estimating continuous functions
from data, offering both predictions and uncertainty estimates. Unlike
parametric models, which assume a fixed functional form, GPR models a
distribution over functions that is updated as new observations are
incorporated [20]. A Gaussian Process (GP) is a collection of random
variables such that any finite subset follows a joint Gaussian
distribution. Formally, it is written as:
\[
   f(x) \sim \mathcal{GP}\bigl(m(x),\, k(x,x')\bigr),
\]
where $m(x)$ is the mean function (often assumed zero) and $k(x,x')$
is the covariance function, or kernel, which encodes the similarity
between inputs. The kernel defines how closely related the outputs
are at different input points. One of the most widely used kernels is
the squared exponential (SE):
\[
   k_{\mathrm{SE}}(x,x') \;=\; \exp\!\left(-\,\frac{\|x - x'\|^{2}}{2l^{2}}\right),
\]
where $l$ is the length scale parameter controlling smoothness (in
our experiments, $l = 2.5$~cm is equal to the radius of the indenter),
and $\|x - x'\|$ is the Euclidean distance between inputs. This kernel
implies smooth, infinitely differentiable functions, making it ideal
for modelling physical quantities such as tissue stiffness.

\vspace{0.3em}

\noindent\textit{\hspace{1em}2)~GPR Training:} Given a set of training
observations $X$ with corresponding outputs $y$, GPR computes a
posterior predictive distribution at new input locations $X^{*}$.
This yields a predicted mean function $\mu^{*}$, representing the
estimated stiffness at unobserved locations, and an associated
covariance $\Sigma^{*}$, which quantifies model uncertainty. As
expected, the uncertainty increases with distance from the observed
samples, reflecting reduced confidence in sparsely explored regions.

\vspace{0.6em}

% ---------- B. Expected Information Density for Elasticity Maps ----------
\noindent\textit{B.~Expected Information Density for Elasticity Maps}

\vspace{0.2em}

\noindent
To prioritise regions for further palpation, we define an EID function
$\xi_{\mathrm{EID}}(x)$ that balances two objectives: exploration of
uncertain areas, and exploitation of predicted high-stiffness regions
and refinement near region boundaries. Formally, we define
\begin{equation}
\tilde{\xi}_{\mathrm{EID}}(x)
  \;=\; (1-\alpha)\!\left[\frac{g(x)}{\int_{X} g(x)}
                         + \frac{\mu(x)}{\int_{X} \mu(x)}\right]
     + \alpha\,\frac{\sigma(x)}{\int_{X} \sigma(x)}
\end{equation}
the un-normalised distribution, where $g(x) = \|\nabla \mu(x)\|$ is
the gradient norm of the predicted elasticity, highlighting
boundaries. The distribution (1) is then normalised as:
\begin{equation}
\xi_{\mathrm{EID}}(x)
  \;=\; \frac{\tilde{\xi}_{\mathrm{EID}}(x)}%
             {\int_{X} \tilde{\xi}_{\mathrm{EID}}(x)\,dx}\,.
\end{equation}
The trade-off between exploration and exploitation is controlled by
the scalar $\alpha \in [0,1]$. Lower values bias the system towards
known high-stiffness or high-gradient areas, while higher values
encourage sampling in regions of high uncertainty. In our
implementation, $\alpha$ varies over time depending on the value of
the ergodic metric in the following way:
\begin{equation}
\alpha(t) \;=\; \frac{\int_{X} e(x,t)\,dx}{\int_{X} e(x,0)\,dx}\,,
\end{equation}
where $e(x,t)$ is the ergodic metric that will be defined in the
following section in (4). According to (3), at the onset of the
exploration, the trajectory is not ergodic and $\alpha(0) = 1$. As a
consequence, the EID in (1) is dominated by the variance of the GPR,
resulting in a purely exploratory behaviour. Regions that have not
yet been sampled exhibit higher uncertainty and are therefore
prioritised for exploration. As informative measurements are
collected, the ergodic metric progressively decreases, leading to a
gradual rebalancing between exploration and exploitation.the spatial
coverage of the trajectory increasingly matches the target EID
distribution, and the exploration naturally shifts towards
exploitation of regions associated with higher estimated stiffness
and pronounced spatial gradients, such as the boundaries of stiff
inclusions. We want to highlight that the EID is computed on-the-fly
and does not exploit prior knowledge about the inspected area.

\vspace{0.6em}

% ---------- C. Trajectory Planning ----------
\noindent\textit{C.~Trajectory Planning}

\vspace{0.2em}

\noindent
The normalised EID $\xi_{\mathrm{EID}}$ is used as the target
distribution for the ergodic planner, which computes a trajectory
that visits each region in proportion to its expected information
content. The HEDAC algorithm [19] is used to generate the trajectory
online.

\vspace{0.3em}

\noindent\textit{\hspace{1em}1)~Ergodic Control:} aims to guide a
robot trajectory to match a target spatial distribution, balancing
exploration and exploitation by spending more time in regions of
higher interest. It models the coverage density using radial basis
functions (RBFs). Given a trajectory
$z : [0,t] \to \Omega \subset \mathbb{R}^{n}$, the normalised
coverage density $c(x,t)$ is:
\[
c(x,t) \;=\; \frac{\tilde{c}(x,t)}{\int_{\Omega} \tilde{c}(x,t)\,dx},
\qquad
\tilde{c}(x,t) \;=\; \frac{1}{t}\int_{0}^{t}
                     \varphi_{\sigma}\!\bigl(x - z(\tau)\bigr)\,d\tau,
\]

% ---------- Bottom license strip ----------
\vfill
\hrule
\vspace{0.3em}
\begin{center}
{\footnotesize\itshape
This work is licensed under a Creative Commons Attribution 4.0
License. For more information, see
\url{https://creativecommons.org/licenses/by/4.0/}\par}
\end{center}

\end{document}
